%\documentclass[11pt,a4paper,twoside]{article}

\usepackage[utf8]{inputenc}

\usepackage[T1,T2A]{fontenc}

\usepackage[english.ancient,russian.ancient]{babel}

\usepackage{graphicx}

\usepackage{array}

\usepackage[doi]{samgtu-bib}

\usepackage{samgtu}

\usepackage{samgtu-name}

\usepackage{mathtools}

\mathtoolsset{showonlyrefs}

\usepackage{desclist}

\usepackage{euscript}

\usepackage{mathrsfs} 

\usepackage{multicol}

\usepackage{mathrsfs}

\usepackage{cite}

\usepackage{cmap}

\usepackage{qrcode}

\allowdisplaybreaks[4]

\usepackage[nodayofweek]{datetime}

\usepackage[thinlines]{easybmat}



\renewcommand{\JournalYear}{2018}

\renewcommand{\JournalVolume}{22}

\renewcommand{\JournalNumber}{4}



\newdate{JournalDateSign}{29}{12}{2018}% Дата подписания Вестника в печать

\renewcommand{\JournalDateSign}{\displaydate{JournalDateSign}}

\newdate{JournalDatePrint}{16}{01}{2019}% Дата выхода Вестника из печати

\renewcommand{\JournalDatePrint}{\displaydate{JournalDatePrint}}

%\renewcommand{\JournalUPL}{16.56} % Усл. печ. л.

%\renewcommand{\JournalUIL}{16.53} % Уч.-изд. л.

\renewcommand{\JournalUPL}{15.24} % Усл. печ. л.

\renewcommand{\JournalUIL}{15.21} % Уч.-изд. л.

\renewcommand{\JournalRegNumber}{220/18} % Регистрационный номер

\renewcommand{\JournalOrderNumber}{2} % Номер заказа



%\renewcommand{\JournalRMonth}{Март} % Месяц выхода Вестника

%\renewcommand{\JournalEMonth}{March} % Месяц выхода Вестника

%\renewcommand{\JournalRMonth}{Июнь} % Месяц выхода Вестника

%\renewcommand{\JournalEMonth}{June} % Месяц выхода Вестника

%\renewcommand{\JournalRMonth}{Сентябрь} % Месяц выхода Вестника

%\renewcommand{\JournalEMonth}{September} % Месяц выхода Вестника

\renewcommand{\JournalRMonth}{Декабрь} % Месяц выхода Вестника

\renewcommand{\JournalEMonth}{December} % Месяц выхода Вестника



\newcommand{\totop}[1]{\raisebox{6pt}[1pt][2pt]{#1}}

\newcommand{\tottop}[1]{\raisebox{12pt}[1pt][2pt]{#1}} 

\newcommand{\totttop}[1]{\raisebox{18pt}[1pt][2pt]{#1}} 

\newcommand{\tottttop}[1]{\raisebox{24pt}[1pt][2pt]{#1}} 

\newcommand{\totttttop}[1]{\raisebox{30pt}[1pt][2pt]{#1}} 

\newcommand{\tobot}[1]{\raisebox{-6pt}[1pt][2pt]{#1}} 

\newcommand{\tobbot}[1]{\raisebox{-12pt}[1pt][2pt]{#1}} 

\newcommand{\tobbbot}[1]{\raisebox{-18pt}[1pt][2pt]{#1}}



\graphicspath{{./00PIC/}}

%\usepackage{samgtu-name}

%

\vsgtu{1577} %registration number

\FirstPage{752}

\LastPage{759}

%

\ShortCommunication

%

%\pdfcrop

%\draft

%

%

%\begin{document}

%

%

%\hfill \qrcode[hyperlink,height=.8in]{http://doi.org/10.14498/vsgtu1577}

%\vspace{-.8in}





\udc{517.956.3}

\msc{35L56}



\rutitle{Задача Коши для системы дифференциальных уравнений гиперболического типа порядка $n$ с~некратными характеристиками}

\rutitleshort{Задача Коши для системы дифференциальных уравнений гиперболического типа\dots}

\rutitlehack{Задача Коши для системы дифференциальных уравнений гиперболического типа порядка $\boldsymbol n$ с~некратными характеристиками}



\entitle{The Cauchy problem for a system of the~hy\-per\-bolic

differential equations of the $n$-th order with the nonmultiple characte\-ristics}

\entitleshort{The Cauchy problem for a system of the hyperbolic differential equations\dots}

\entitlehack{The Cauchy problem for a system of the hyperbolic

differential equations of the $\boldsymbol n$-th order with~the~nonmultiple~characteristics}



\ruauthor{А.~А.~Андреев\affil{1}, Ю.~О.~Яковлева\affil{2}}{А\,н\,д\,р\,е\,е\,в~А.~А.,\ Я\,к\,о\,в\,л\,е\,в\,а~Ю.~О.}

\enauthor{A.~A.~Andreev\affil{1}, J.~O.~Yakovleva\affil{2}}{A\,n\,d\,r\,e\,e\,v~A.~A.,\ Y\,a\,k\,o\,v\,l\,e\,v\,a~J.~O.}



\ruaffil{\rusamgtu.\and \rusgau.}

\enaffil{\ensamgtu.\and \ensgau.}







\ruauthorsdetails{2}{% Количество авторов

\fullauthor{\rupid{38541}{Александр Анатольевич Андреев}}

%\CAuthor % Автор-корреспондент

\orcid{0000-0002-6611-6685} % ORCID ID автора 

\regalia{кандидат физико-математических наук, доцент} 

\position{доцент} 

\dept{каф. прикладной математики и информатики} 

\email{andre01071948@yandex.ru}

\fullauthor{\rupid{126662}{Юлия Олеговна Яковлева}}

\CAuthor % Автор-корреспондент

\orcid{0000-0002-9839-3740} % ORCID ID автора 

\regalia{кандидат физико-математических наук} 

\position{доцент} 

\dept{каф. математики и бизнес"=информатики} 

\email{julia.yakovleva@mail.ru}

}



\enauthorsdetails{2}{

\fullauthor{\enpid{38541}{Aleksandr A. Andreev}}

\orcid{0000-0002-6611-6685} % ORCID ID автора 

\regalia{Cand.  Phys. \& Math. Sci., Professor} 

\position{Associate Professor} 

\dept{Dept. of Applied Mathematics \& Computer Science}

\email{andre01071948@yandex.ru}

\fullauthor{\enpid{126662}{Julia O. Yakovleva}}

\CAuthor % Автор-корреспондент

\orcid{0000-0002-9839-3740} % ORCID ID автора 

\regalia{Cand. Phys. \& Math. Sci., Associate Professor} 

\position{Associate Professor} 

\dept{Dept. of Mathematics \& Business Informatics}

\email{julia.yakovleva@mail.ru}

}



\rukeywords{гиперболическое дифференциальное уравнение порядка $n$,

cистема уравнений гиперболического типа порядка $n$,  некратные

характеристики, метод общих решений, задача Коши, формула

Даламбера}



\enkeywords{$n$-th order hyperbolic differential equation, system

of the hyperbolic differential equations of the $n$-th order,

nonmultiple characteristics, method of the general solutions,

Cauchy problem,  D'Alembert formula}





\ruabstract{Рассмотрена задача Коши для дифференциального

гиперболического уравнения порядка $n$ с~некратными

характеристиками. Приведено регулярное решение задачи Коши для

дифференциального уравнения гиперболического типа порядка $n$ с~некратными характеристиками. 

Получено решение задачи Коши для системы дифференциальных уравнений гиперболического типа

порядка $n$, не содержащей производных меньше порядка $n$,  с~некратными характеристиками в~случае коммутирующих матричных коэффициентов. Как результат исследований сформулирована теорема о~существовании и единственности регулярного решения задачи Коши для системы дифференциальных  уравнений гиперболического типа порядка $n$ с~некратными характеристиками.}



\enabstract{In the paper the Cauchy problem is considered for the hyperbolic differential equation of the $n$-th order with the nonmultiple characteristics. The regular solution of the Cauchy problem for the hyperbolic differential equation of the $n$-th order with the nonmultiple characteristics is considered. In the paper the solution of the Cauchy problem for the system of the hyperbolic differential equations of the $n$-th order with the nonmultiple characteristics is considered. The existence and uniqueness theorem for the regular solution of the Cauchy problem for the system of the hyperbolic differential equations of the $n$-th order with the nonmultiple characteristics is considered as the result of the research.}



\newdate{andreyakoa}{17}{11}{2017}

\date{\displaydate{andreyakoa}}

\newdate{andreyakob}{13}{12}{2017}

\revisiondate{\displaydate{andreyakob}}

\newdate{andreyakoc}{18}{12}{2017}

\accepteddate{\displaydate{andreyakoc}}



\newdate{andreyakoe}{25}{12}{2017}

\renewcommand{\JournalDataOnline}{\displaydate{andreyakoe}}



\vspace{5mm}



%\ForPeerReview

\makerutitle



\newpage







\Section[n]{Введение} Исследование краевых задач для

дифференциальных уравнений гиперболического типа и~систем

гиперболических уравнений  является одним из важных разделов

современной теории дифференциальных уравнений с частными

производными. Интерес многих исследователей к данной проблеме

объясняется как теоретической важностью получаемых результатов,

так и их значимыми практическими

приложениями~[\citen{yakovleva:bib:Ali,yakovleva:bib:Bicadze,

yakovleva:bib:Tamotu,yakovleva:bib:Popivanov}]. 

Краевые задачи, в~том числе и задачу Коши, для уравнений гиперболического типа и~систем уравнений  гиперболического типа с некратными

характеристиками в ряде случаев можно решить не только методом

Римана~[\citen{yakovleva:bib:Rieman,yakovleva:bib:Gegalov}], но и

методом общих решений~[\citen{yakovleva:bib:Mus}].



\smallskip



\Section[n]{Предварительные сведения} В

работе~[\citen{yakovleva:bib:Yakovleva10}] рассмотрена и решена

задача Коши для дифференциального уравнения гиперболического типа

 порядка $n$ общего вида с некратными характеристиками

 \begin{equation}

\sum\limits_{k=0}^{n}a_{k}\dfrac{\partial^{n}}{\partial x^{n-k}\partial y^{k}}u(x,\,y)=0,\label{eq:Yako1}

\end{equation}

где $a_{k} $ "--- действительные ненулевые постоянные.



Уравнение~\eqref{eq:Yako1} является строго

гиперболическим по Петровскому~[\citen{yakovleva:bib:Petrovskiy}],

то есть имеет $n$ некратных характеристик.

Тогда характеристическое уравнение для уравнения~\eqref{eq:Yako1}

$$

a_{0}\lambda^{n}+a_{1}\lambda^{n-1}+\ldots+a_{n-1}\lambda+a_{n}=0

$$

 имеет

$n$ различных действительных отличных от нуля корней

$\lambda_{1},\,\lambda_{2},\,\dots,\,\lambda_{n}$ таких, что

$$\lambda_{1}+\lambda_{2}+\dots+\lambda_{n}=-\frac{a_{1}}{a_{0}},\quad 

\lambda_{1}\lambda_{2}+\lambda_{1}\lambda_{3}+\dots+\lambda_{n-1}\lambda_{n}=\frac{a_{2}}{a_{0}},$$

$$\lambda_{1}\lambda_{2}\lambda_{3}+\lambda_{1}\lambda_{2}\lambda_{4}+\dots+\lambda_{n-2}\lambda_{n-1}\lambda_{n}=-\frac{a_{3}}{a_{0}},\quad \dots,$$

$$\lambda_{1}\lambda_{2}\cdots\lambda_{n-1}+\lambda_{2}\lambda_{3}\cdots\lambda_{n}+\dots+\lambda_{1}\lambda_{3}\cdots\lambda_{n}=(-1)^{n-1}\frac{a_{n-1}}{a_{0}},$$

$$\lambda_{1}\lambda_{2}\cdots\lambda_{n}=(-1)^{n}\frac{a_{n}}{a_{0}}.$$



Известно~[\citen{yakovleva:bib:Korzuk,yakovleva:bib:Yakovleva1,yakovleva:bib:Yakovleva11}], что

общее решение уравнения~\eqref{eq:Yako1} принадлежит классу $n$

раз непрерывно дифференцируемых функций

$C^{n}(\mathbb{R}\times\mathbb{R})$ и может быть представлено в

виде

$$

u(x,y)=f_{1}(y+\lambda_{1}x)+f_{2}(y+\lambda_{2}x)+\ldots

+f_{n}(y+\lambda_{n}x).

$$



\smallskip



{\small\sc Задача Коши.} {\it Найти регулярное решение

$u\left(x,\,y\right)\in C^{n}(\mathbb{R}\times\mathbb{R})$

уравнения~\eqref{eq:Yako1} в плоскости независимых переменных

$D=\{\left(x,\,y\right): x \in \mathbb{R},\, y \in \mathbb{R}\},$

которое на линии $y=0$ удовлетворяет условиям

 \begin{equation}

\dfrac{\partial^{k} u(x,\,y)}{\partial l^{k}}\Bigr|_{y=0}=\tau_{n-k}(x),\quad k=0, 1, n-1,\label{eq:Yako2}

\end{equation}

%\begin{gather}

%u\left(x,\,y\right)|_{y=0}=\tau_{n}(x),\notag

%\\\dfrac{\partial u}{\partial n}|_{y=0}=\tau_{n-1}(x),\notag

%\\\dfrac{\partial^{2}u}{\partial n^{2}}|_{y=0}=\tau_{n-2}(x),\notag

%\\\ldots \notag

%\\\dfrac{\partial^{n-1}u}{\partial n^{n-1}}|_{y=0}=\tau_{1}(x),\label{eq:Yako2}

%\end{gather}

где $\tau_{1}(x),\,\tau_{2}(x),\,\ldots,\tau_{n}(x) \in

C^{n}(\mathbb{R}),$ ${l}$ "--- нормаль к нехарактеристической линии $y=0.$

}

\smallskip



Функция $u\left(x,\,y\right)\in

C^{n}(\mathbb{R}\times\mathbb{R})$, имеющая в плоскости $D$

все

непрерывные частные производные, которые входят в

уравнение~\eqref{eq:Yako1}, и удовлетворяющая

уравнению~\eqref{eq:Yako1} и условиям задачи Коши~\eqref{eq:Yako2}

в обычном смысле, называется регулярным

решением 

задачи Коши~\eqref{eq:Yako1}, \eqref{eq:Yako2} в плоскости независимых переменных~[\citen{yakovleva:bib:Yakovleva3}].



Решением задачи~\eqref{eq:Yako1}, \eqref{eq:Yako2} является функция

\begin{equation}

u(x,y)=\sum\limits_{k=1}^{n}f_{k}(y+\lambda_{k}x),\label{eq:Yako3}

\end{equation}

где

\begin{multline}

f_{k}(\lambda_{k}x)=\lambda_{k}^{n-1}\frac{(-1)^{n+k+1}}{(n-2)!}\frac{a_{n}}{a_{0}}\int

_{0}^{x}\tau_{1}(s)(x-s)^{n-2}ds+ 

\\

+\lambda_{k}^{n-1}\frac{(-1)^{n+k+1}}{(n-3)!}\frac{a_{n-1}}{a_{0}}\int _{0}^{x}\tau_{2}(s)(x-s)^{n-3}ds+\dots+\lambda_{k}^{n-1}(-1)^{n+k+1}\tau_{n}(x).

\label{eq:Yako4}

\end{multline}



Явная формула~\eqref{eq:Yako3} регулярного решения рассмотренной

задачи Коши~[\citen{yakovleva:bib:Yakovleva10}] является

естественным обобщением известной формулы

Даламбера~[\citen{yakovleva:bib:Tihonov}].



Справедлива 



\smallskip



{\small\sc Теорема~1.} {\it 

Существует единственное регулярное решение задачи Коши}~\eqref{eq:Yako1}, \eqref{eq:Yako2} {\it

 $u\left(x,\,y\right)\in C^{n}(\mathbb{R}\times\mathbb{R}),$  имеющее

вид~\eqref{eq:Yako3} при условии$,$ что функции $\tau_{1}(x),$ $\tau_{2}(x),\ldots,\tau_{n}(x) \in

C^{n}(\mathbb{R}).$

}



\smallskip



\Section[n]{Задача Коши для системы дифференциальных уравнений

гиперболического типа, не содержащей производных меньше порядка~$\boldsymbol n$}

Рассмотрим систему дифференциальных уравнений гиперболического типа порядка

$n$ с двумя независимыми переменными $x, y$ на

плоскости $D$, не содержащую производных порядка меньше $n$,

\begin{equation}

\sum\limits_{k=0}^{n}A_{k}\dfrac{\partial^{n}}{\partial x^{n-k}\partial y^{k}}U(x,\,y)=0,\label{eq:Yako5}

\end{equation}

где $U(x,y)=(u^{1}(x,y),u^{2}(x,y))^{\top}$ "--- двумерная вектор"=функция, 

$A_{k}$ "--- постоянные $(2\times2)$"=матрицы.



Будем предполагать, что матрицы $A_{0},\,A_{1},\ldots, A_{n}$

попарно коммутируют~[\citen{yakovleva:bib:Bellman}].

Следовательно, существует такое преобразование подобия $T$,

которое одновременно приводит матрицы $A_{0},\,A_{1},\ldots,

A_{n}$ к диагональной форме

$$

T^{-1}A_{0}T=\Lambda_{A_{0}},\quad 

T^{-1}A_{1}T=\Lambda_{A_{1}},\quad 

\ldots,\quad 

T^{-1}A_{n}T=\Lambda_{A_{n}}.

$$



Если матрицы $A_{0},\,A_{1},\ldots,A_{n}$ коммутирующие, то

матрицы $\Lambda_{A_{0}},$ $\Lambda_{A_{1}},$ $\ldots,$ $\Lambda_{A_{n}}$,

полученные преобразованием подобия~[\citen{yakovleva:bib:Gantmaher}], также попарно коммутируют.



Пусть матрицы

$\Lambda_{A_{0}},$ $\Lambda_{A_{1}},$ $\ldots,$ $\Lambda_{A_{n}}$ таковы,

что имеют различные ненулевые действительные собственные значения.



\smallskip



{\small\sc Задача Коши.} {\it  Найти регулярное решение

$U\left(x,\,y\right)\in C^{n}(\mathbb{R}\times\mathbb{R})$ системы

дифференциальных уравнений~\eqref{eq:Yako5} в плоскости $D,$

которое на линии $y=0$ удовлетворяет условиям

\begin{equation}

\dfrac{\partial^{k} U(x,\,y)}{\partial l^{k}}\Bigr|_{y=0}=S_{n-k}(x),\quad k= 0, 1, n-1 ,\label{eq:Yako6}

\end{equation}

%\begin{gather}

%U\left(x,\,0\right)=S_{n}(x),\, \dfrac{\partial U}{\partial

%n}\left(x,0\right)=S_{n-1}(x),\,\ldots,\dfrac{\partial^{n-1}

%U}{\partial n^{n-1}}\left(x,0\right)=S_{1}(x)\label{eq:Yako6}

%\end{gather}

где $S_{1}(x),\,S_{2}(x),\ldots,S_{n}(x)\in C^{n}(\mathbb{R})$ "---

заданные вектор"=функции$,$

${l}$ "--- нормаль к нехарактеристической линии $y=0.$}



\smallskip



Рассмотрим замену

$$U=TV,\quad V(x,y)=(v^{1}(x,y),v^{2}(x,y))^{\top}$$ при $\det T\neq

0$, которая позволит совершить переход к системе вида

\begin{equation}

\Lambda_{A_{0}}V_{xxx\ldots x}+\Lambda_{A_{1}}V_{xx\ldots xy}+\ldots+\Lambda_{A_{n-1}}V_{n-1}u_{x\ldots

yyy}+\Lambda_{A_{n}}V_{n}u_{yyy\ldots y}=0.\label{eq:Yako7}

\end{equation}



Таким образом, исследуемая система разделена на отдельные

уравнения:

$$

\begin{array}{l}

a_{0}^{1}v^{1}_{xxx\ldots x}+a_{1}^{1}v^{1}_{xx\ldots xy}+\ldots+a_{n-1}^{1}v^{1}_{x\ldots yyy}+a_{n}^{1}v^{1}_{yyy\ldots y}=0, \\[2mm]

a_{0}^{2}v^{2}_{xxx\ldots x}+a_{1}^{2}v^{2}_{xx\ldots

xy}+\ldots+a_{n-1}^{2}v^{2}_{x\ldots

yyy}+a_{n}^{2}v^{2}_{yyy\ldots y}=0.

\end{array} 

$$



Каждое характеристическое уравнение системы~\eqref{eq:Yako7} имеет $n$ различных ненулевых корней

$\lambda_{1},\,\lambda_{2},\ldots,\lambda_{n},\,\mu_{1},\,\mu_{2},\ldots,\mu_{n}$

соответственно.



Применяя приведенные выше исследования, для каждого уравнения

указанной системы  решение задачи Коши может быть получено в

регулярном виде:

$$

v^{1}(x,y)=\sum\limits_{k=1}^{n}f^{1}_{k}(y+\lambda_{k}x), \quad 

v^{2}(x,y)=\sum\limits_{k=1}^{n}f^{2}_{k}(y+\mu_{k}x),

$$

где $f^{1}_{k},\,f^{2}_{k}$ определяются по

формуле~\eqref{eq:Yako4}.



Решение матричного уравнения $U=TV$ является решением  задачи

Коши~\eqref{eq:Yako5}, \eqref{eq:Yako6}.



Таким образом, существование решения задачи Коши~\eqref{eq:Yako5}, \eqref{eq:Yako6}

доказано конструктивно. 

Существование и единственность решения задачи Коши для линейной системы уравнений гиперболического типа с~аналитическими коэффициентами доказано Хольмгреном~[\citen{yakovleva:bib:Holmgren}].







Приведенные исследования позволяют сформулировать следующую

теорему.



\smallskip

{\small\sc Теорема 2.} {\it В плоскости $D$

существует единственное регулярное решение $U\left(x,\,y\right)\in

C^{n}(\mathbb{R}\times\mathbb{R})$ задачи Коши {\rm \eqref{eq:Yako5}, \eqref{eq:Yako6}}

при условии$,$ что вектор"=функции $S_{1}(x),\,S_{2}(x),\ldots,S_{n}(x)\in C^{n}(\mathbb{R}).$

}

\smallskip



Полученный результат является обобщением результатов, полученных авторами в~работах~[\citen{yakovleva:bib:Yakovleva10, yakovleva:bib:Yakovleva1, yakovleva:bib:Yakovleva11}].





\AdditionalContent{Конкурирующие интересы}{Заявляем, что в отношении авторства и публикации

этой статьи конфликта интересов не имеем.} 

\AdditionalContent{Авторский вклад и ответственность}{Все  авторы  принимали  участие  в  разработке  концепции  \mbox{статьи}  и  в  написании  рукописи. Авторы несут  полную  ответственность  за  предоставление  окончательной рукописи в печать. Окончательная  версия рукописи была одобрена всеми авторами.} 

\AdditionalContent{Финансирование}{Исследование выполнялось без финансирования.}





\begin{thebibliography}{99}



\Bibitem{yakovleva:bib:Ali}

\by Ali Raeisian~S.~M.

\paper Effective Solution of Riemann Problem for Fifth Order Improperly Elliptic Equation on a Rectangle

\jour AJCM

\vol 2 

\issue 4

\yr 2012

\pages 282--286

\crossref{10.4236/ajcm.2012.24038}



\RBibitem{yakovleva:bib:Bicadze}  

\by Бицадзе~А.~В.

\book Уравнения математической физики

\yr 1982

\publ Наука

\publaddr М.

\totalpages 336

\mathscinet{http://www.ams.org/mathscinet-getitem?mr=701245}

\zmath{https://zbmath.org/?q=an:03882793}

%\misctransl

%\by Bitsadze~A.~V.

%\book Uravneniia matematicheskoi fiziki \rm [Equations of mathematical physics]

%\yr 1982

%\publ Nauka

%\publaddr Moscow

%\lang In Russian

%\totalpages 336







\Bibitem{yakovleva:bib:Tamotu} 

\paper Gevrey Wellposedness of the Cauchy Problem for the Hyperbolic Equations of Third Order with Coefficients Depending Only on Time

\jour  Publications of the Research Institute for Mathematical Sciences

\by Kinoshita~T.

\vol 34 

\yr 1998

\issue 3 

\pages 249--270

\crossref{10.2977/prims/1195144695}





\Bibitem{yakovleva:bib:Popivanov}  

\by Nikolov~A., Popivanov~N.

\paper Singular solutions to Protter's problem for (3+1)-D degenerate wave equation

\procinfo 8–13 June 2012; Sozopol, Bulgaria

\serial AIP Conf. Proc.

\yr 2012

\vol 1497

\pages 233--238

\crossref{10.1063/1.4766790}

\scopus{http://www.scopus.com/record/display.url?origin=inward&eid=2-s2.0-84873199179}



\Bibitem{yakovleva:bib:Rieman}

\by Rieman~B.

\paper Ueber die Fortpflanzung ebener Luftwellen von endlicher Schwingungsweite (Aus dem achten Bande der Abhandlungen der K\"oniglichen Gesellschaft der Wissenschaften zu G\"ottingen. 1860.)

\inbook Bernard Riemann's Gesammelte mathematische Werke und wissenschaftlicher Nachlass

\yr 2009

\pages 145--164

\eds R.~Dedekind, H.~M.~Weber

\publ BiblioLife

\publaddr United States

\lang In German

\crossref{10.1017/cbo9781139568050.009}



\RBibitem{yakovleva:bib:Gegalov} 

\by Жегалов~В.~И., Миронов~А.~Н. 

\book Дифференциальные уравнения со старшими частными производными 

\publaddr Казань

\publ Казанское математическое общество 

\yr 2001 

\totalpages 226

\zmath{https://zbmath.org/?q=an:1038.35500}



\RBibitem{yakovleva:bib:Mus} 

\by Мусхелишвили~Н.~И.

\book Некоторые основные задачи математической теории упругости 

\publaddr М.

\publ Наука

\yr 1966

\totalpages   707

\mathscinet{http://www.ams.org/mathscinet-getitem?mr=202367}

\zmath{https://zbmath.org/?q=an:0151.36201}



\RBibitem{yakovleva:bib:Yakovleva10}  

\by Андреев~А.~А., Яковлева~Ю.~О. 

\paper Задача Коши для уравнения гиперболического типа порядка $n$ общего вида

с~некратными характеристиками

\jour Вестн. Сам. гос. техн. ун-та. Сер. Физ.-мат. науки

\yr 2016

\vol 20

\issue 2

\pages 241--248

\mathnet{http://mi.mathnet.ru/vsgtu1490}

\crossref{10.14498/vsgtu1490}

\elib{http://elibrary.ru/item.asp?id=27126223}



\RBibitem{yakovleva:bib:Petrovskiy}  

\by Петровский~И.~Г.

\book Избранные труды. Системы уравнений с~частными производными. Алгебраическая геометрия

\yr 1986

\publ Наука

\publaddr М.

\totalpages 500

\mathscinet{http://www.ams.org/mathscinet-getitem?mr=871873}

\zmath{https://zbmath.org/?q=an:0603.01018}

%\misctransl

%\by Petrovsky~I.~G.

%\book Izbrannye trudy. Sistemy uravnenii s~chastnymi proizvodnymi. Algebraicheskaia geometriia \rm [Selected works. Systems of partial differential equations. Algebraic geometry]

%\yr 1986

%\publ Nauka

%\publaddr Moscow

%\totalpages 504





\RBibitem{yakovleva:bib:Korzuk}  

\by Корзюк~В.~И., Чеб~Е.~С., Ле~Тхи~Тху

\paper Решение смешанной задачи для биволнового уравнения методом характеристик

\jour Тр. Ин-та матем.

\yr 2010

\vol 18

\issue 2

\pages 36--54

\mathnet{http://mi.mathnet.ru/timb16}

\zmath{https://zbmath.org/?q=an:05863489}



\RBibitem{yakovleva:bib:Yakovleva1} 

\by Яковлева~Ю.~О.

\paper Аналог формулы Даламбера для гиперболического уравнения третьего порядка с~некратными характеристиками

\jour Вестн. Сам. гос. техн. ун-та. Сер. Физ.-мат. науки

\yr 2012

\issue 1(26)

\pages 247--250

\mathnet{http://mi.mathnet.ru/vsgtu1028}

\crossref{10.14498/vsgtu1028}



\RBibitem{yakovleva:bib:Yakovleva11}  

\by Андреев~А.~А., Яковлева~Ю.~О. 

\paper Задача Коши для~системы уравнений гиперболического типа четвертого порядка

общего~вида с~некратными характеристиками

\jour Вестн. Сам. гос. техн. ун-та. Сер. Физ.-мат. науки

\yr 2014

\issue 4(37)

\pages 7--15

\mathnet{http://mi.mathnet.ru/vsgtu1349}

\crossref{10.14498/vsgtu1349}

\elib{http://elibrary.ru/item.asp?id=23464547}



\RBibitem{yakovleva:bib:Yakovleva3} 

\by Андреев~А.~А., Яковлева~Ю.~О. 

\by А.~А.~Андреев, Ю.~О.~Яковлева

\paper Характеристическая задача для одного гиперболического дифференциального уравнения третьего порядка с некратными характеристиками

\jour Изв. Сарат. ун-та. Нов. сер. Сер. Математика. Механика. Информатика

\yr 2013

\vol 13

\issue 1(2)

\pages 3--6

\mathnet{http://mi.mathnet.ru/isu361}



\RBibitem{yakovleva:bib:Tihonov} 

\by Тихонов~А.~Н., Самарский~А.~А.

\book Уравнения математической физики

\yr 1972

\publ Наука

\publaddr М.

\totalpages 736

\mathscinet{http://www.ams.org/mathscinet-getitem?mr=344638}

\zmath{https://zbmath.org/?q=an:03417939}

%\misctransl

%\by Tikhonov~A.~N., Samarskii~A.~A.

%\book Uravneniia matematicheskoi fiziki \rm [Equations of Mathematical Physics]

%\yr 1972

%\publ Nauka

%\publaddr Moscow

%\lang In Russian

%\totalpages 736





\Bibitem{yakovleva:bib:Bellman} 

\by Bellman~R.

\book Introduction to matrix analysis

\bookinfo 2nd ed., Reprint of the 1970 Orig.

\serial Classics in Applied Mathematics

\yr 1997

\vol 19

\publ Society for Industrial and Applied Mathematics

\publaddr Philadelphia, PA

\totalpages xxviii+403

\mathscinet{http://www.ams.org/mathscinet-getitem?mr=1455129}

\zmath{https://zbmath.org/?q=an:01034200}





\RBibitem{yakovleva:bib:Gantmaher} 

\by Гантмахер~Ф.~Р.

\book Теория матриц

\yr 1988

\publ Наука

\publaddr М.

\mathscinet{http://www.ams.org/mathscinet-getitem?mr=986246}

\zmath{https://zbmath.org/?q=an:0666.15002}

\totalpages 549



\Bibitem{yakovleva:bib:Holmgren}  

\by Holmgren~E.

\paper Sur les syst\`emes lin\'eaires aux d\'eriv\'ees partielles du premier ordre deux variables ind\'ependantes \`a caract\'eristiques r\'eelles et distinetes

\jour Arkiv f. Mat., Astr. och Fys.

\yr 1909

\vol 5

\issue 1

\lang In Swedish

\totalpages 13

\zmath{https://zbmath.org/?q=an:40.0415.01}



\end{thebibliography}







\newpage







\makeentitle



\AdditionalContent{Competing interests}{We declare that we have no conflicts of interests with the authorship and publication of this article.} 

\AdditionalContent{Authors' contributions and responsibilities}{Each author has participated in the article concept development and in the manuscript writing. The authors are absolutely responsible for submitting the final manuscript in print. Each author has approved the final version of manuscript.} 



\AdditionalContent{Funding}{The research has not had any funding.}







\begin{thebibliography}{99}



\Bibitem{yakovleva:bib:Ali:en}

\by Ali Raeisian~S.~M.

\paper Effective Solution of Riemann Problem for Fifth Order Improperly Elliptic Equation on a Rectangle

\jour AJCM

\vol 2 

\issue 4

\yr 2012

\pages 282--286

\crossref{10.4236/ajcm.2012.24038}



\Bibitem{yakovleva:bib:Bicadze:en}  

\by Bitsadze~A.~V.

\book Uravneniia matematicheskoi fiziki \rm [Equations of mathematical physics]

\yr 1982

\publ Nauka

\publaddr Moscow

\lang In Russian

\totalpages 336







\Bibitem{yakovleva:bib:Tamotu:en} 

\paper Gevrey Wellposedness of the Cauchy Problem for the Hyperbolic Equations of Third Order with Coefficients Depending Only on Time

\jour  Publications of the Research Institute for Mathematical Sciences

\by Kinoshita~T.

\vol 34 

\yr 1998

\issue 3 

\pages 249--270

\crossref{10.2977/prims/1195144695}





\Bibitem{yakovleva:bib:Popivanov:en}  

\by Nikolov~A., Popivanov~N.

\paper Singular solutions to Protter's problem for (3+1)-D degenerate wave equation

\procinfo 8–13 June 2012; Sozopol, Bulgaria

\serial AIP Conf. Proc.

\yr 2012

\vol 1497

\pages 233--238

\crossref{10.1063/1.4766790}

\scopus{http://www.scopus.com/record/display.url?origin=inward&eid=2-s2.0-84873199179}



\Bibitem{yakovleva:bib:Rieman:en}

\by Rieman~B.

\paper Ueber die Fortpflanzung ebener Luftwellen von endlicher Schwingungsweite (Aus dem achten Bande der Abhandlungen der K\"oniglichen Gesellschaft der Wissenschaften zu G\"ottingen. 1860.)

\inbook Bernard Riemann's Gesammelte mathematische Werke und wissenschaftlicher Nachlass

\yr 2009

\pages 145--164

\eds R.~Dedekind, H.~M.~Weber

\publ BiblioLife

\publaddr United States

\lang In German

\crossref{10.1017/cbo9781139568050.009}



\Bibitem{yakovleva:bib:Gegalov:en} 

\lang In Russian

\by Zhegalov~V.~I., Mironov~A.~N. 

\book Differentsial'nye uravneniia so starshimi chastnymi proizvodnymi \rm [Differential equations with highest partial derivatives]

\publaddr Kazan'

\publ Kazanskoe matematicheskoe obshchestvo 

\yr 2001 

\totalpages 226

\zmath{https://zbmath.org/?q=an:1038.35500}



\Bibitem{yakovleva:bib:Mus:en} 

\lang In Russian

\by Muskhelishvili~N.~I.

\book Nekotorye osnovnye zadachi matematicheskoi teorii uprugosti  \rm [Some basic problems of the mathematical theory of elasticity]

\publaddr Moscow

\publ Nauka

\yr 1966

\totalpages   707





\Bibitem{yakovleva:bib:Yakovleva10:en}  

\lang In Russian

\by Andreev~A.~A., Yakovleva~J.~O.

\paper The Cauchy problem for a general hyperbolic differential equation of the $n$-th order

with~the~nonmultiple characteristics

\jour Vestn. Samar. Gos. Tekhn. Univ., Ser. Fiz.-Mat. Nauki \rm [J. Samara State Tech. Univ., Ser. Phys. Math. Sci.]

\yr 2016

\vol 20

\issue 2

\pages 241--248

\crossref{10.14498/vsgtu1490}



\Bibitem{yakovleva:bib:Petrovskiy:en}  

\by Petrovsky~I.~G.

\book Izbrannye trudy. Sistemy uravnenii s~chastnymi proizvodnymi. Algebraicheskaia geometriia \rm [Selected works. Systems of partial differential equations. Algebraic geometry]

\yr 1986

\publ Nauka

\publaddr Moscow

\totalpages 504

\lang In Russian





\Bibitem{yakovleva:bib:Korzuk:en}  

\lang In Russian

\by Korzyuk~V.~I., Cheb~E.~S., Le~Thi~Thu

\paper Solution of the mixed problem for the biwave equation by the method of characteristics

\jour Tr. Inst. Mat.

\yr 2010

\vol 18

\issue 2

\pages 36--54



\Bibitem{yakovleva:bib:Yakovleva1:en} 

\lang In Russian

\by Yakovleva~J.~O.

\paper The analogue of D'Alembert formula for hyperbolic differential equation of the third order with nonmultiple characteristics

\jour Vestn. Samar. Gos. Tekhn. Univ., Ser. Fiz.-Mat. Nauki \rm [J. Samara State Tech. Univ., Ser. Phys. Math. Sci.]

\yr 2012

\issue 1(26)

\pages 247--250

\crossref{10.14498/vsgtu1028}



\Bibitem{yakovleva:bib:Yakovleva11:en}  

\lang In Russian

\by Andreev~A.~A., Yakovleva~J.~O.

\paper Cauchy Problem For the System Of the General Hyperbolic Differential Equations

Of the Forth Order With Nonmultiple Characteristics

\jour Vestn. Samar. Gos. Tekhn. Univ., Ser. Fiz.-Mat. Nauki \rm [J. Samara State Tech. Univ., Ser. Phys. Math. Sci.]

\yr 2014

\issue 4(37)

\pages 7--15

\crossref{10.14498/vsgtu1349}





\Bibitem{yakovleva:bib:Yakovleva3:en} 

\lang In Russian

\by Andreev~A.~A., Yakovleva~J.~O.

\paper The Characteristic Problem for one Hyperbolic Differentional Equation of the Third Order with Nonmultiple Characteristics

\jour Izv. Saratov Univ. (N.S.), Ser. Math. Mech. Inform.

\yr 2013

\vol 13

\issue 1(2)

\pages 3--6



\Bibitem{yakovleva:bib:Tihonov:en} 

\by Tikhonov~A.~N., Samarskii~A.~A.

\book Uravneniia matematicheskoi fiziki \rm [Equations of Mathematical Physics]

\yr 1972

\publ Nauka

\publaddr Moscow

\lang In Russian

\totalpages 736





\Bibitem{yakovleva:bib:Bellman:en} 

\by Bellman~R.

\book Introduction to matrix analysis

\bookinfo 2nd ed., Reprint of the 1970 Orig.

\serial Classics in Applied Mathematics

\yr 1997

\vol 19

\publ Society for Industrial and Applied Mathematics

\publaddr Philadelphia, PA

\totalpages xxviii+403

\mathscinet{http://www.ams.org/mathscinet-getitem?mr=1455129}

\zmath{https://zbmath.org/?q=an:01034200}





\Bibitem{yakovleva:bib:Gantmaher:en} 

\lang In Russian

\by Gantmakher~F.~R.

\book Teoriia matrits \rm [Theory of matrices]

\yr 1988

\publ Nauka

\publaddr Moscow

\totalpages 549





\Bibitem{yakovleva:bib:Holmgren:en}  

\by Holmgren~E.

\paper Sur les syst\`emes lin\'eaires aux d\'eriv\'ees partielles du premier ordre deux variables ind\'ependantes \`a caract\'eristiques r\'eelles et distinetes

\jour Arkiv f. Mat., Astr. och Fys.

\yr 1909

\vol 5

\issue 1

\lang In Swedish

\totalpages 13

\zmath{https://zbmath.org/?q=an:40.0415.01}



\end{thebibliography}







%\end{document}

